\chapter*{ }\addcontentsline{toc}{chapter}{Лицензия и дополнительное соглашение}
\markboth{ЛИЦЕНЗИЯ И ДОПОЛНИТЕЛЬНОЕ СОГЛАШЕНИЕ}{ЛИЦЕНЗИЯ И ДОПОЛНИТЕЛЬНОЕ СОГЛАШЕНИЕ}
\markright{}
\label{chapter:license}
\begin{center}\Large
Публичная лицензия Creative Commons ``Указание авторства --- Некоммерческое использование --- Без производных произведений'' 4.0 Международная (Creative Commons Attribution-NonCommercial-NoDerivatives 4.0 International Public License)
\end{center}

\small

Осуществляя Лицензированные права (Licensed Rights) (определённые ниже), Вы принимаете и соглашаетесь соблюдать условия настоящей Публичной лицензии Creative Commons ``Указание авторства --- Некоммерческое использование --- Без производных произведений'' 4.0 Международная (Creative Commons Attribution-NonCommercial-NoDerivatives 4.0 International Public License) (далее --- ``Публичная лицензия'' (Public License)). Насколько настоящая Публичная лицензия (Public License) может быть истолкована как договор, Вам предоставляются Лицензированные права (Licensed Rights) в обмен на Ваше принятие данных условий, а Лицензиар (Licensor) предоставляет Вам такие права в обмен на выгоды, которые Лицензиар (Licensor) получает от предоставления Лицензированного материала (Licensed Material) на условиях настоящей Публичной лицензии (Public License).

\subsection*{Раздел 1 -- Определения.}
\begin{itemize}
\item[a.] \textbf{Адаптированный материал (Adapted Material)} означает материал, охраняемый Авторским правом и аналогичными правами (Copyright and Similar Rights), который является производным от Лицензированного материала (Licensed Material) или основан на нём и в котором Лицензированный материал (Licensed Material) переведён, изменён, адаптирован, трансформирован или иным образом модифицирован способом, требующим разрешения по Авторскому праву и аналогичным правам (Copyright and Similar Rights), принадлежащим Лицензиару (Licensor). В целях настоящей Публичной лицензии (Public License), если Лицензированный материал (Licensed Material) является музыкальным произведением, исполнением или фонограммой, Адаптированный материал (Adapted Material) всегда создаётся, когда Лицензированный материал (Licensed Material) синхронизируется во временном соотношении с движущимся изображением.

\item[b.] \textbf{Авторское право и аналогичные права (Copyright and Similar Rights)} означают авторское право и/или аналогичные права, тесно связанные с авторским правом, включая, помимо прочего, права на исполнение, передачу в эфир, фонограмму и Суи генерис права на базы данных (Sui Generis Database Rights), независимо от того, как обозначены или классифицированы эти права. В целях настоящей Публичной лицензии (Public License) права, указанные в разделе \href{https://creativecommons.org/licenses/by-nc-nd/4.0/legalcode#s2b}{2(b)(1)-(2)}, не являются Авторским правом и аналогичными правами (Copyright and Similar Rights).

\item[c.] \textbf{Эффективные технологические меры (Effective Technological Measures)} означают такие меры, которые при отсутствии надлежащих полномочий не могут быть обойдены в соответствии с законами, выполняющими обязательства по статье 11 Договора ВОИС об авторском праве, принятого 20 декабря 1996 года, и/или аналогичным международным соглашениям.

\item[d.] \textbf{Исключения и ограничения (Exceptions and Limitations)} означают добросовестное использование (fair use), добросовестную сделку (fair dealing) и/или любое иное исключение или ограничение Авторского права и аналогичных прав (Copyright and Similar Rights), применимое к Вашему использованию Лицензированного материала (Licensed Material).

\item[e.] \textbf{Лицензированный материал (Licensed Material)} означает художественное или литературное произведение, базу данных или иной материал, к которому Лицензиар (Licensor) применил настоящую Публичную лицензию (Public License).

\item[f.] \textbf{Лицензированные права (Licensed Rights)} означают права, предоставленные Вам на условиях настоящей Публичной лицензии (Public License), которые ограничены всеми Авторскими правами и аналогичными правами (Copyright and Similar Rights), применимыми к Вашему использованию Лицензированного материала (Licensed Material) и которые Лицензиар (Licensor) имеет право лицензировать.

\item[g.] \textbf{Лицензиар (Licensor)} означает физическое(ие) или юридическое(ие) лицо(а), предоставляющее права на условиях настоящей Публичной лицензии (Public License).

\item[h.] \textbf{Некоммерческий (NonCommercial)} означает, что использование не направлено в первую очередь на получение коммерческой выгоды или денежного вознаграждения. В целях настоящей Публичной лицензии (Public License) обмен Лицензированного материала (Licensed Material) на иной материал, охраняемый Авторским правом и аналогичными правами (Copyright and Similar Rights), путём цифрового обмена файлами или аналогичными средствами, является Некоммерческим (NonCommercial), при условии отсутствия выплаты денежного вознаграждения в связи с таким обменом.

\item[i.] \textbf{Распространение (Share)} означает предоставление материала публике любым способом или процессом, требующим разрешения по Лицензированным правам (Licensed Rights), таким как воспроизведение, публичная демонстрация, публичное исполнение, распространение, доведение до всеобщего сведения, сообщение или импортирование, а также обеспечение доступа материала для публики, включая способы, при которых члены публики могут получить доступ к материалу из места и в время, индивидуально выбранные ими.

\item[j.] \textbf{Суи генерис права на базы данных (Sui Generis Database Rights)} означают права, отличные от авторского права, возникающие из Директивы 96/9/EC Европейского парламента и Совета от 11 марта 1996 года о правовой охране баз данных, с поправками и/или последующими редакциями, а также других существенно эквивалентных прав в любой точке мира.

\item[k.] \textbf{Вы (You)} означает физическое или юридическое лицо, осуществляющее Лицензированные права (Licensed Rights) на условиях настоящей Публичной лицензии (Public License). Ваш (Your) имеет соответствующее значение.
\end{itemize}

\subsection*{Раздел 2 -- Область действия.}
\begin{itemize}
\item[a.] \textbf{Предоставление лицензии.}
	\begin{enumerate}
	\item При условии соблюдения условий настоящей Публичной лицензии (Public License) Лицензиар (Licensor) настоящим предоставляет Вам всемирную, безвозмездную, несублицензируемую, неисключительную, безотзывную лицензию на осуществление Лицензированных прав (Licensed Rights) в отношении Лицензированного материала (Licensed Material) для:
	\begin{enumerate}
		\item[A.] воспроизведения и Распространения (Share) Лицензированного материала (Licensed Material), полностью или частично, исключительно для Некоммерческих (NonCommercial) целей; и
		\item[B.] создания и воспроизведения, но не Распространения (Share), Адаптированного материала (Adapted Material) исключительно для Некоммерческих (NonCommercial) целей.
	\end{enumerate}
	\item \underline{Исключения и ограничения (Exceptions and Limitations).} Для устранения сомнений, в случаях, когда Исключения и ограничения (Exceptions and Limitations) применимы к Вашему использованию, настоящая Публичная лицензия (Public License) не применяется, и Вам не требуется соблюдать её условия.
	\item \underline{Срок действия.} Срок действия настоящей Публичной лицензии (Public License) указан в разделе \href{https://creativecommons.org/licenses/by-nc-nd/4.0/legalcode#s6a}{6(a)}.
	\item \underline{Носители и форматы; разрешённые технические изменения.} Лицензиар (Licensor) уполномочивает Вас осуществлять Лицензированные права (Licensed Rights) на всех носителях и в любых форматах, известных в настоящее время или созданных в будущем, а также вносить технические изменения, необходимые для этого. Лицензиар (Licensor) отказывается от права и/или соглашается не заявлять о каком-либо праве или полномочии запрещать Вам вносить технические изменения, необходимые для осуществления Лицензированных прав (Licensed Rights), включая технические изменения, необходимые для обхода Эффективных технологических мер (Effective Technological Measures). В целях настоящей Публичной лицензии (Public License) простое внесение изменений, разрешённых настоящим разделом \href{https://creativecommons.org/licenses/by-nc-nd/4.0/legalcode#s2a4}{2(a)(4)}, никогда не создаёт Адаптированный материал (Adapted Material).
	\item \underline{Последующие получатели.}
		\begin{enumerate}
		\item[A.] \underline{Предложение от Лицензиара (Licensor) --- Лицензированный материал (Licensed Material).} Каждый получатель Лицензированного материала (Licensed Material) автоматически получает предложение от Лицензиара (Licensor) осуществлять Лицензированные права (Licensed Rights) на условиях настоящей Публичной лицензии (Public License).
		\item[B.] \underline{Отсутствие ограничений для последующих получателей.} Вы не можете предлагать или устанавливать какие-либо дополнительные или иные условия для Лицензированного материала (Licensed Material) или применять к нему какие-либо Эффективные технологические меры (Effective Technological Measures), если это ограничивает осуществление Лицензированных прав (Licensed Rights) любым получателем Лицензированного материала (Licensed Material).
		\end{enumerate}
\item \underline{Отсутствие одобрения.} Ничто в настоящей Публичной лицензии (Public License) не составляет и не может быть истолковано как разрешение утверждать или подразумевать, что Вы или Ваше использование Лицензированного материала (Licensed Material) связаны с Лицензиаром (Licensor) или иными лицами, которым должно быть указано авторство в соответствии с разделом \href{https://creativecommons.org/licenses/by-nc-nd/4.0/legalcode#s3a1Ai}{3(a)(1)(A)(i)}, или что Лицензиар (Licensor) или такие лица спонсируют, одобряют или предоставляют официальный статус Вам или Вашему использованию.
	\end{enumerate}
\item[b.] \textbf{Иные права.}
	\begin{enumerate}
	\item Моральные права, такие как право на неприкосновенность произведения, не лицензируются на условиях настоящей Публичной лицензии (Public License), равно как и права на публичность, неприкосновенность частной жизни и/или другие аналогичные персональные права; однако, насколько это возможно, Лицензиар (Licensor) отказывается от таких прав и/или соглашается не заявлять о любых таких правах, принадлежащих ему, в ограниченной степени, необходимой для позволения Вам осуществлять Лицензированные права (Licensed Rights), но не иначе.
	\item Патентные и товарные права не лицензируются на условиях настоящей Публичной лицензии (Public License).
	\item Насколько это возможно, Лицензиар (Licensor) отказывается от права взимать с Вас авторские вознаграждения за осуществление Лицензированных прав (Licensed Rights), независимо от того, напрямую или через организацию по коллективному управлению правами на основании любой добровольной или отказной законодательной или обязательной лицензионной схемы. Во всех остальных случаях Лицензиар (Licensor) прямо сохраняет за собой право взимать такие авторские вознаграждения, включая случаи, когда Лицензированный материал (Licensed Material) используется не для Некоммерческих (NonCommercial) целей.
	\end{enumerate}
\end{itemize}

\subsection*{Раздел 3 -- Условия лицензии.}
Ваше осуществление Лицензированных прав (Licensed Rights) прямо подчиняется следующим условиям.
\begin{itemize}
\item[a.] \textbf{Указание авторства (Attribution).}
	\begin{enumerate}
	\item Если Вы Распространяете (Share) Лицензированный материал (Licensed Material), Вы должны:
		\begin{enumerate}
		\item[A.] сохранить следующее, если оно предоставлено Лицензиаром (Licensor) вместе с Лицензированным материалом (Licensed Material):
			\begin{enumerate}
			\item идентификацию создателя(ей) Лицензированного материала (Licensed Material) и любых других лиц, которым должно быть указано авторство, любым разумным способом, требуемым Лицензиаром (Licensor) (включая псевдоним, если указан);
			\item уведомление об авторском праве;
			\item уведомление, ссылающееся на настоящую Публичную лицензию (Public License);
			\item уведомление, ссылающееся на отказ от гарантий;
			\item URI или гиперссылку на Лицензированный материал (Licensed Material) в мере разумной возможности;
			\end{enumerate}
		\item[B.] указать, изменяли ли Вы Лицензированный материал (Licensed Material), и сохранить указание о любых предыдущих изменениях; и
		\item[C.] указать, что Лицензированный материал (Licensed Material) лицензирован на условиях настоящей Публичной лицензии (Public License), и включить текст или URI или гиперссылку на настоящую Публичную лицензию (Public License).
		\end{enumerate}
Для устранения сомнений, Вам не разрешено на основании настоящей Публичной лицензии (Public License) Распространять (Share) Адаптированный материал (Adapted Material).
\item Вы можете удовлетворить условиям раздела \href{https://creativecommons.org/licenses/by-nc-nd/4.0/legalcode#s3a1A}{3(a)(1)} любым разумным способом в зависимости от носителя, средств и контекста, в котором Вы Распространяете (Share) Лицензированный материал (Licensed Material). Например, может быть разумно удовлетворить условиям путём предоставления URI или гиперссылки на ресурс, содержащий требуемую информацию.
\item Если Лицензиар (Licensor) потребует, Вы должны удалить любую информацию, требуемую разделом \href{https://creativecommons.org/licenses/by-nc-nd/4.0/legalcode#s3a1A}{3(a)(1)(A)}, в мере разумной возможности.
	\end{enumerate}
\end{itemize}

\subsection*{Раздел 4 -- Суи генерис права на базы данных (Sui Generis Database Rights).}
Если Лицензированные права (Licensed Rights) включают Суи генерис права на базы данных (Sui Generis Database Rights), применимые к Вашему использованию Лицензированного материала (Licensed Material):
\begin{itemize}
\item[a.]
для устранения сомнений, раздел \href{https://creativecommons.org/licenses/by-nc-nd/4.0/legalcode#s2a1}{2(a)(1)} предоставляет Вам право извлекать, повторно использовать, воспроизводить и Распространять (Share) всё или существенную часть содержимого базы данных исключительно для Некоммерческих (NonCommercial) целей при условии, что Вы не Распространяете (Share) Адаптированный материал (Adapted Material);
\item[b.] если Вы включаете всё или существенную часть содержимого базы данных в базу данных, в отношении которой у Вас имеются Суи генерис права на базы данных (Sui Generis Database Rights), то такая база данных, в отношении которой у Вас имеются Суи генерис права на базы данных (Sui Generis Database Rights) (но не её отдельные элементы), является Адаптированным материалом (Adapted Material); и
\item[c.] Вы должны соблюдать условия раздела \href{https://creativecommons.org/licenses/by-nc-nd/4.0/legalcode#s3a}{3(a)}, если Вы Распространяете (Share) всё или существенную часть содержимого базы данных.
\end{itemize}
Для устранения сомнений, настоящий раздел \href{https://creativecommons.org/licenses/by-nc-nd/4.0/legalcode#s4}{4} дополняет и не заменяет Ваши обязательства по настоящей Публичной лицензии (Public License) в случаях, когда Лицензированные права (Licensed Rights) включают иные Авторские права и аналогичные права (Copyright and Similar Rights).

\subsection*{Раздел 5 -- Отказ от гарантий (Disclaimer of Warranties) и ограничение ответственности (Limitation of Liability).}
\begin{itemize}
\item[\textbf{a.}] \textbf{Если иное не взято на себя Лицензиаром (Licensor) отдельно, насколько это возможно, Лицензиар (Licensor) предоставляет Лицензированный материал (Licensed Material) ``как есть'' и ``по мере доступности'', и не даёт никаких заверений или гарантий любого рода в отношении Лицензированного материала (Licensed Material), явных, подразумеваемых, предусмотренных законом или иных. Это включает, помимо прочего, гарантии права собственности, товарной пригодности, пригодности для конкретной цели, отсутствия нарушений, отсутствия скрытых или иных дефектов, точности, наличия или отсутствия ошибок, известных или обнаружимых. В случаях, когда отказ от гарантий не допускается полностью или частично, данный отказ может не применяться к Вам.}
\item[\textbf{b.}] \textbf{Насколько это возможно, ни при каких обстоятельствах Лицензиар (Licensor) не несёт перед Вами ответственности по любой правовой теории (включая, помимо прочего, неосторожность) или иным образом за любые прямые, специальные, косвенные, случайные, последующие, штрафные, примерные (exemplary) или иные убытки, издержки, расходы или ущерб, возникшие в связи с настоящей Публичной лицензией (Public License) или использованием Лицензированного материала (Licensed Material), даже если Лицензиар (Licensor) был предупреждён о возможности таких убытков, издержек, расходов или ущерба. В случаях, когда ограничение ответственности не допускается полностью или частично, данное ограничение может не применяться к Вам.}
\item[c.] Отказ от гарантий и ограничение ответственности, указанные выше, должны толковаться таким образом, чтобы насколько возможно наиболее близко приблизиться к абсолютному отказу от ответственности и полному освобождению от ответственности.
\end{itemize}

\subsection*{Раздел 6 -- Срок действия и прекращение.}
\begin{itemize}
\item[a.] Настоящая Публичная лицензия (Public License) действует в течение срока Авторского права и аналогичных прав (Copyright and Similar Rights), лицензируемых настоящей лицензией. Однако если Вы нарушаете условия настоящей Публичной лицензии (Public License), Ваши права по настоящей Публичной лицензии (Public License) автоматически прекращаются.
\item[b.] Если Ваше право на использование Лицензированного материала (Licensed Material) прекращено в соответствии с разделом \href{https://creativecommons.org/licenses/by-nc-nd/4.0/legalcode#s6a}{6(a)}, оно восстанавливается:
	\begin{enumerate}
	\item автоматически с даты устранения нарушения, при условии, что оно устранено в течение 30 дней с момента Вашего обнаружения нарушения; или
	\item по явному восстановлению Лицензиаром (Licensor).
	\end{enumerate}
Для устранения сомнений, настоящий раздел \href{https://creativecommons.org/licenses/by-nc-nd/4.0/legalcode#s6b}{6(b)} не затрагивает никакого права Лицензиара (Licensor) на взыскание средств за Ваши нарушения настоящей Публичной лицензии (Public License).
\item[c.] Для устранения сомнений, Лицензиар (Licensor) может также предлагать Лицензированный материал (Licensed Material) на отдельных условиях или прекратить его распространение в любое время; однако это не прекращает действие настоящей Публичной лицензии (Public License).
\item[d.] Разделы \href{https://creativecommons.org/licenses/by-nc-nd/4.0/legalcode#s1}{1}, \href{https://creativecommons.org/licenses/by-nc-nd/4.0/legalcode#s5}{5}, \href{https://creativecommons.org/licenses/by-nc-nd/4.0/legalcode#s6}{6}, \href{https://creativecommons.org/licenses/by-nc-nd/4.0/legalcode#s7}{7} и \href{https://creativecommons.org/licenses/by-nc-nd/4.0/legalcode#s8}{8} сохраняют действие после прекращения настоящей Публичной лицензии (Public License).
\end{itemize}

\subsection*{Раздел 7 -- Прочие условия.}
\begin{itemize}
\item[a.] Лицензиар (Licensor) не связан никакими дополнительными или иными условиями, сообщёнными Вами, если только не согласовано явно.
\item[b.] Любые договорённости, договоренности или соглашения в отношении Лицензированного материала (Licensed Material), не указанные в настоящем документе, являются отдельными от и независимыми от условий настоящей Публичной лицензии (Public License).
\end{itemize}

\subsection*{Раздел 8 -- Толкование.}
\begin{itemize}
\item[a.] Для устранения сомнений, настоящая Публичная лицензия (Public License) не уменьшает, не ограничивает, не сужает и не устанавливает условий для любого использования Лицензированного материала (Licensed Material), которое могло бы быть законно осуществлено без разрешения по настоящей Публичной лицензии (Public License).
\item[b.] Насколько это возможно, если какое-либо положение настоящей Публичной лицензии (Public License) признаётся неисполнимым, оно должно быть автоматически реформировано в минимально необходимой степени для обеспечения его исполнимости. Если положение не может быть реформировано, оно должно быть исключено из настоящей Публичной лицензии (Public License) без ущерба для исполнимости остальных условий.
\item[c.] Ни одно условие настоящей Публичной лицензии (Public License) не может быть отказано, и ни одно несоблюдение не может быть санкционировано, если только Лицензиар (Licensor) не согласился на это явно.
\item[d.] Ничто в настоящей Публичной лицензии (Public License) не составляет и не может быть истолковано как ограничение или отказ от каких-либо привилегий и иммунитетов, применимых к Лицензиару (Licensor) или Вам (You), включая в отношении судебных процессов любой юрисдикции или органа власти.
\end{itemize}

Creative Commons не является стороной своих публичных лицензий. Несмотря на это, Creative Commons может применить одну из своих публичных лицензий к материалу, который она публикует, и в таких случаях будет считаться ``Лицензиаром'' (Licensor). Текст публичных лицензий Creative Commons передан в общественное достояние на условиях \href{https://creativecommons.org/publicdomain/zero/1.0/legalcode}{\textsl{посвящения общественному достоянию CC0}} (CC0 Public Domain Dedication). За исключением ограниченной цели указания, что материал распространяется на условиях публичной лицензии Creative Commons, или в ином случае, разрешённом политиками Creative Commons, опубликованными на \href{https://creativecommons.org/policies}{creativecommons.org/policies}, Creative Commons не уполномочивает использование товарного знака ``Creative Commons'' или любого другого товарного знака или логотипа Creative Commons без её предварительного письменного согласия, включая, помимо прочего, в связи с любыми несанкционированными изменениями любой из её публичных лицензий или любыми иными договорённостями, соглашениями или договорами об использовании лицензированного материала. Для устранения сомнений, настоящий абзац не является частью публичных лицензий.

С Creative Commons можно связаться по адресу \href{creativecommons.org}{creativecommons.org}.

\chapter*{ }
\label{chapter:contrib}

\begin{center}\Large 
Дополнительное соглашение (Sideletter) по внесению вкладов в Руководство MoonMath по zk-SNARK (далее --- ``\textbf{Sideletter}'')
\end{center}\label{sideletter}

между Least Authority TFA GmbH, Thaerstraße 28a, 10249 Berlin (далее именуемой ``\textbf{Least Authority}'') и любым физическим или юридическим лицом, представляющим вклады в Руководство MoonMath (далее именуемым ``\textbf{Вы}'' или ``\textbf{Ваш}''). 
\begin{center}\bfseries
Преамбула
\end{center}
\begin{itemize}
\item[(A)] Least Authority является первоначальным создателем так называемого Руководства MoonMath по zk-SNARK (далее --- ``\textbf{Руководство}'' (Manual)), которое служит ресурсом для всех, кто заинтересован в понимании и раскрытии потенциала так называемой технологии ``zk-SNARK'' (``\textbf{zk-SNARK}''). Аббревиатура zk-SNARK расшифровывается как ``Zero-Knowledge Succinct Non-Interactive Argument of Knowledge'' и относится к криптографической технике, при которой можно доказать владение определённой информацией, не раскрывая саму информацию. Большинство объяснений с трудом проясняют, как и почему они работают. Ресурсы разбросаны по блогам и библиотекам Github. Это приводит к высокому барьеру для входа, тем самым замедляя широкое распространение zk-SNARK и связанных с ними технологий повышения конфиденциальности. 
\item[(B)] Least Authority хочет изменить это с помощью Руководства, продолжая развивать его как основанный на сообществе проект для сбора полезной и практической информации о zk-SNARK. Авторы третьих сторон, такие как Вы, должны иметь возможность вносить части, идеи и практическую информацию в Руководство. 
\item[(C)] Само Руководство лицензировано на условиях Публичной лицензии Creative Commons, версия \textsl{Attri\-bution-NonCommercial-NoDerivatives 4.0 International} (``\textbf{CC BY-NC-ND-4.0}''), которая допускает использование и распространение, а также изменение Руководства. Однако если Вы изменяете Руководство или создаёте ``\textbf{Адаптированный материал}'' (Adapted Material) Руководства в смысле раздела 1.a. лицензии CC BY-NC-ND-4.0, Вам не разрешается распространять его, поскольку раздел 3.a.1. пункт 2 лицензии CC BY-NC-ND-4.0 запрещает распространение Адаптированного материала (Adapted Material). 
\item[(D)] Если Вы желаете участвовать в Руководстве, Вы можете представить Адаптированный материал (Adapted Material) по Руководству, а также материал, созданный независимо от Руководства (``\textbf{Независимые произведения}'' (Independent Creations)), в Least Authority. Если Вы заинтересованы в добавлении значительного Вклада (Contribution) в Руководство, пожалуйста, свяжитесь с Least Authority напрямую по адресу \href{mailto:mmm@leastauthority.com}{mmm@leastauthority.com}, и мы можем обсудить, может ли Ваш вклад быть рассмотрен индивидуально на иных условиях.
\item[(E)] Предметом настоящего Дополнительного соглашения (Sideletter) является лицензирование Вашего Вклада (Contribution) в пользу Least Authority. 
\end{itemize}
Стороны договорились о нижеследующем:

\begin{center}
$\S$1\\\bfseries
Лицензия на представленный Вами Вклад (Contribution)
\end{center}

\begin{itemize}
\item[(1)] Вы можете вносить любые письменные произведения, графические материалы, методы расчёта, компиляции информации, базы данных или любые другие произведения авторства, включая любые изменения или дополнения к Руководству, созданные Вами, представляя их в Least Authority для включения в Руководство, независимо от того, являются ли они Независимыми произведениями (Independent Creations) или Адаптированным материалом (Adapted Material) (каждое из них --- ``\textbf{Вклад}'' (Contribution)). ``\textbf{Представление}'' (Submission) в этом смысле включает любую форму электронной, устной или письменной коммуникации, направленной в Least Authority по адресу \href{mailto:mmm@leastauthority.com}{mmm@leastauthority.com} или размещённой на https://github.com/LeastAuthority/moonmath-manual. Для ясности: Least Authority не обязана включать Ваш Вклад (Contribution) в Руководство.

\item[(2)] Настоящим Вы предоставляете Least Authority бессрочное, всемирное, неисключительное, сублицензируемое, безотзывное и безвозмездное право использовать, изменять, редактировать, обнародовать и распространять Ваш Вклад (Contribution) в материальной и нематериальной форме или любым иным способом, известным в настоящее время или разработанным в будущем, в первоначальном или изменённом виде (в пределах запрета порчи), а также комбинировать его в первоначальном или изменённом виде с Руководством или в составе Руководства (далее --- ``\textbf{Лицензия}'' (License)). Лицензия (License) включает как минимум все права, необходимые для лицензирования Вклада (Contribution) на условиях CC BY-NC-ND-4.0, и в частности позволяет Least Authority использовать, изменять, редактировать, обнародовать и распространять в материальной и нематериальной форме или любым иным способом, известным в настоящее время или разработанным в будущем, Вклад (Contribution) как часть Руководства. Least Authority настоящим принимает предоставление Лицензии (License). 
\item[(3)] Если Least Authority решит, что Ваш Вклад (Contribution) или его части должны быть включены в Руководство, Least Authority обеспечит следующее: 
	\begin{itemize}
	\item[a)] Вклад (Contribution) как часть Руководства лицензируется на условиях CC BY-NC-ND-4.0, 
	\item[b)] Вы будете указаны как Автор вклада (Contributor) (включая псевдоним, если указан) в Руководстве. 
	\end{itemize}
Правило § 1 (3) b) применяется только в том случае, если Ваше имя или псевдоним указаны вместе с Вкладом (Contribution). 
\item[(4)] В случае если Least Authority решит, что только части или редакции Вашего Вклада (Contribution) будут включены в Руководство, Least Authority уведомит Вас в разумный срок и получит Ваше согласие на лицензирование частей / редакций Вашего Вклада (Contribution) в соответствии с §1 (2). Согласие не требуется, если Least Authority вносит только редакторские правки. В случае если Вы представляете Ваш Вклад (Contribution) без контактной информации, положение § 1 (4) не применяется, поскольку у Least Authority нет возможности получить Ваше согласие.      
\item[(5)] В случае если Least Authority решит, что Ваш Вклад (Contribution) не будет частью Руководства, Least Authority примет разумные меры для информирования Вас о своём решении в разумный срок после Вашего Представления (Submission). Предоставленная Вами Лицензия (License) на Вклад (Contribution) прекращается с момента принятия Least Authority решения не включать Вклад (Contribution) в Руководство. 
\end{itemize}

\begin{center}
$\S$2\\\bfseries
Отказ от гарантий (Disclaimer)
\end{center}
\begin{itemize}
\item[(1)] Если иное не взято на себя Вами отдельно, насколько это возможно, Вы предоставляете Вклад (Contribution) ``как есть'' и ``по мере доступности'', и не даёте никаких заверений или гарантий любого рода в отношении Вклада (Contribution), явных, подразумеваемых, предусмотренных законом или иных. Это включает, помимо прочего, гарантии права собственности, товарной пригодности, пригодности для конкретной цели, отсутствия нарушений, отсутствия скрытых или иных дефектов, точности, наличия или отсутствия ошибок, известных или обнаружимых. В случаях, когда отказ от гарантий не допускается полностью или частично, данный отказ может не применяться к Вам.
\item[(2)] Насколько это возможно, ни при каких обстоятельствах Вы не несёте перед нами ответственности по любой правовой теории (включая, помимо прочего, неосторожность) или иным образом за любые прямые, специальные, косвенные, случайные, последующие, штрафные, примерные (exemplary) или иные убытки, издержки, расходы или ущерб, возникшие в связи с настоящим Дополнительным соглашением (Sideletter) или использованием Вклада (Contribution), даже если Вас предупреждали о возможности таких убытков, издержек, расходов или ущерба. В случаях, когда ограничение ответственности не допускается полностью или частично, данное ограничение может не применяться к Вам.
\item[(3)] Отказ от гарантий и ограничение ответственности, указанные выше, должны толковаться таким образом, чтобы насколько возможно наиболее близко приблизиться к абсолютному отказу от ответственности и полному освобождению от ответственности.
\end{itemize}

\begin{center}
$\S$3\\\bfseries
Разное (Miscellaneous)
\end{center}

\begin{itemize}
\item[(1)] Настоящее Дополнительное соглашение (Sideletter) действительно без подписи. Оно заключается между Вами и Least Authority в момент представления Вами Вклада (Contribution) в Least Authority. 
\item[(2)] Настоящее Дополнительное соглашение (Sideletter), его толкование и любые внедоговорные обязательства в связи с ним регулируются материальным правом Германии. Конвенция ООН о договорах международной купли-продажи товаров (CISG) не применяется.
\item[(3)] Английские термины, используемые в настоящем Дополнительном соглашении (Sideletter), описывают исключительно немецкие правовые концепции и не должны толковаться в соответствии со значением, придаваемым им в какой-либо иной юрисдикции, кроме Германии. В случае вставки немецкого термина в скобках и/или курсивом только он (а не связанный с ним английский термин) является авторитетным для целей толкования соответствующего термина при любом его использовании в настоящем Соглашении.
\item[(4)] Если одно или несколько положений настоящего Дополнительного соглашения (Sideletter) являются или становятся недействительными или неисполнимыми полностью или частично, это не затрагивает действительность и исполнимость остальных положений настоящего Дополнительного соглашения (Sideletter). Вместо любых недействительных или не включённых в настоящее Дополнительное соглашение (Sideletter) Общих условий бизнеса (\textsl{Allgemeine Geschäftsbedingungen}) применяются законодательные положения (§ 306 (2) Германского гражданского уложения (BGB)). Во всех остальных случаях стороны договорятся о действительном положении для замены недействительного или неисполнимого положения, которое наиболее точно отражает первоначальную экономическую цель, при условии, что дополнительное толкование Дополнительного соглашения (\textsl{ergänzende Vertragsauslegung}) не имеет приоритета или невозможно.
\item[(5)] Изменения и дополнения к настоящему Дополнительному соглашению (Sideletter) действительны только в письменной форме. Это также относится к любому изменению данного положения о письменной форме. 
\item[(6)] Любые споры, возникающие из или в связи с настоящим Дополнительным соглашением (Sideletter), включая споры о его заключении, обязательности, изменении и прекращении, подлежат рассмотрению исключительно компетентным судом в Берлине, Германия, если это юридически возможно.
\end{itemize}
